\section{Обсуждение}

\textbf{Что подтвердилось.} Открытословарная карта почти никогда не накрывает истинный след объекта целиком, поэтому без калиброванного запаса планировщик систематически «срезает углы» возле запретных объектов. Калибровка только в пространстве меток (покадровая CP) у открытого словаря вырождается в «запретно всё занятое» и при этом ломается от сантиметров дрейфа: она требует, чтобы каждая клетка объекта попала в множество, а дрейф сдвигает клетки. Радиус, откалиброванный только по меткам, держится, пока дрейф мал по сравнению с ошибкой меток (до $\sim$10~см ATE в наших данных), и проваливается при одометрическом дрейфе. Радиус только по позе недопокрывает при любом реалистичном дрейфе, потому что не видит ошибок детектора. Совместная калибровка расстояния промаха держит номинальное покрытие на всех уровнях и даёт наименьший из валидных запасов.

\textbf{Что оказалось не так, как ожидалось.} Эффект композиции проявился не как недопокрытие раздельной схемы (ср.~\cite{baral}), а как её переплата радиусом на $36$--$38\,\%$ (на L4): синтетический дрейф независим от ошибок детектора, и сумма двух квантилей получается с запасом. Недопокрытие раздельной схемы следует ожидать при положительной корреляции ошибок (например, когда в тёмных и бедных текстурой местах одновременно растут ошибка позы и путаница меток) --- это проверяемая гипотеза для реального SLAM.

\textbf{Главная находка для лаборатории.} Узкое место --- не дрейф и не пограничные ошибки сегментации, а уверенная подмена класса целого объекта. Для тумбы под ТВ это происходит в 2 сценах из 21, т.~е. чаще, чем допускает $\alpha = 0{,}1$, и честная калибровка отказывается давать гарантию. Метрики качества карты (mIoU, f-mIoU) этого не отражают: они усредняют по клеткам и классам, а безопасность определяется наихудшим объектом в сцене.

\textbf{Ограничения и угрозы валидности.}
\begin{itemize}[nosep]
  \item Дрейф синтетический (одометрическая модель без замыканий циклов); реальные ошибки SLAM коррелированы с видом сцены.
  \item Одно условие освещения: остальные условия OSMa-Bench не опубликованы. Ось освещения --- в дальнейших шагах.
  \item Сцены ReplicaCAD --- перестановки мебели одной квартиры; обменяемость сцен выполняется, но разнообразие ограничено. Классы и порог $\lambda_0$ выбраны по dev-сцене, но набор запретных классов для миссий (без тумбы) уточнён после анализа покрытия --- это явно оговорено.
  \item Гарантия маргинальна по сценам; при 13 калибровочных сценах распределение покрытия по разбиениям широкое (10-й перцентиль 0{,}75--0{,}83 у совместной ветки при среднем 0{,}93), что согласуется с бета-распределением покрытия split-CP при малом $n$.
  \item Восприятие --- детектор рамок с масками по глубине, а не полный картограф (ConceptGraphs, BBQ); вывод о подмене классов нужно перепроверить на них.
  \item 2D-карта и плоское движение; ошибка локализации во время исполнения плана не моделируется.
\end{itemize}
